\begin{figure}[t]
\centering\includegraphics[width=\columnwidth]{lane_sweep.pdf}
\caption{同一预训练权重派生 H32/N640/K512 输入下，运算通道数对核心周期数的影响。全部得分和选择结果一致。该核心独立测试不包含封装层的位置排序及原生 DMA／存储等待；结果输出周期取决于测试平台的 ready 序列。不同实现之间不预设时钟频率或 PPA 等价。}
\label{fig:lanes}
\end{figure}
在相同输入及相同握手调度下，8／16／32 路实现分别需要 608,365、444,525 和 362,605 个周期。各实现均计算 2,621,440 个乘积项，并选择相同的 512 个索引。从 8 路增加至 16 路、从 16 路增加至 32 路时，优化前后周期数之比分别为 1.369$\times$ 和 1.226$\times$。增加运算通道可缩短整数乘积的分批计算时间，但缩放转换与 FP32 归约仍串行执行。装载、选择和结果输出的开销也未消除，因此该结果不意味着性能线性扩展，亦不意味着不同实现具有相同的可达频率。

表~\ref{tab:core} 中 Top-512 用例的 444,514 个周期来自另一组随机背压序列，与受控通道数实验中 16 路实现的 444,525 个周期不属于完全相同的时序实验。

\begin{table}[t]
\centering\caption{16 路评分／堆核心在 TSMC 28 nm HPC+ 标准单元库下使用 Design Compiler 获得的实际综合映射结果。数组由触发器实现；未完成布局布线、寄生互连提取、门级等价验证或物理签核。}
\label{tab:synthesis}
\small
\begin{tabular}{lr}\toprule
指标 & 测量值／配置值\\\midrule
库工艺角 & TT 0.9 V 25 $^{\circ}$C\\
目标周期（ns） & 2.000\\
映射单元面积（$\mu$m$^2$） & 301,340.93\\
叶单元数 & 250,241\\
时序单元数 & 53,993\\
建立时间裕量（ns） & +0.000006\\
保持时间裕量（ns） & +0.005131\\
电气约束违例条目数 & 0\\
约束覆盖检查 & 通过\\
\bottomrule\end{tabular}
\end{table}
在上述映射核心模型中，报告的建立时间与保持时间检查均满足所给约束。但建立时间裕量仅为 6 fs，保持时间裕量为 5.131 ps；接近零的名义建立时间裕量不足以证明物理实现具有稳健的工作频率。电气约束违例条目统计转换时间、电容、扇出及其他报告的电气限制，同一对象若出现在不同类别中则分别计数。该检查既不同于建立／保持时间检查，也不同于版图 DRC。独立的时钟与 I/O 约束覆盖检查通过，其中复位信号按文档规定设为假路径。组合单元面积为 163,838.47 $\mu$m$^2$，非组合单元面积为 137,502.45 $\mu$m$^2$，占总面积的 45.6\%。后者同时包含控制状态与触发器存储，不能单独归因于 Q 或堆存储。全部 1,142 个 Liberty 单元的面积均与同名 LEF 单元的宽高乘积精确一致，工艺库发布版本的兼容关系亦已核对；实际 DC 数据库中的单元面积也逐项精确匹配。LEF SIZE 的尺寸单位为微米~\cite{lefdef}，因此这里得到的是标准单元占用面积，而非布局后的核心面积或芯片面积。综合对象不包含 AXI 封装层、DMA、处理器模型及 UCIe 系统。综合使用理想时钟，未建模 CTS 与实际时钟网络负载；对于报告含 53,993 个负载的时钟网，工具采用 1,000 的扇出估计值。独立的已保存设计回读核查了映射与时序报告，但这不构成等价性证明。仿真时钟仅为测试设置，不代表已经实现的芯片工作频率。
